Both inputs must have the same secret key, ring, modulus $q$, and scale $\Delta=q/t$, with $t\mid q$ for the stated plaintext arithmetic modulo $t$. Throughout these operation chapters, $\textsf{GLWE}_{S,\sigma}(\text{phase})$ is shorthand for a ciphertext with that phase; an evaluated ciphertext need not have the same fresh Gaussian distribution or standard deviation $\sigma$. All noise bounds use integer lifts before coefficient reduction.

\textbf{- Reference:} 
\href{https://www.zama.ai/post/tfhe-deep-dive-part-2}{TFHE Deep Dive - Part II - Encodings and linear leveled operations}~\cite{tfhe-2}


$ $

Suppose we have two GLWE ciphertexts encrypting two different plaintexts $M^{\langle 1 \rangle}, M^{\langle 2 \rangle}$:

$\textsf{GLWE}_{S, \sigma}(\Delta M^{\langle 1 \rangle} + E^{\langle 1 \rangle}) = \textsf{ct}^{\langle 1 \rangle} = ( A_0^{\langle 1 \rangle}, A_1^{\langle 1 \rangle}, \ldots, A_{k-1}^{\langle 1 \rangle}, B^{\langle 1 \rangle}) \in \mathcal{R}_{\langle n,q \rangle}^{k + 1}$

$\textsf{GLWE}_{S, \sigma}(\Delta M^{\langle 2 \rangle} + E^{\langle 2 \rangle}) = \textsf{ct}^{\langle 2 \rangle} = ( A_0^{\langle 2 \rangle}, A_1^{\langle 2 \rangle}, \ldots, A_{k-1}^{\langle 2 \rangle}, B^{\langle 2 \rangle}) \in \mathcal{R}_{\langle n,q \rangle}^{k + 1}$

$ $

\noindent Let's define the following ciphertext addition operation: 

$\textsf{ct}^{\langle 1 \rangle} + \textsf{ct}^{\langle 2 \rangle} = ( A_0^{\langle 1 \rangle} + A_0^{\langle 2 \rangle}, \text{ } A_1^{\langle 1 \rangle} + A_1^{\langle 2 \rangle}, \ldots, A_{k-1}^{\langle 1 \rangle} + A_{k-1}^{\langle 2 \rangle}, \text{ } B^{\langle 1 \rangle} + B^{\langle 2 \rangle} )$

$ $

\noindent Then, the following is true:

\begin{tcolorbox}[title={\textbf{\tboxlabel{\ref*{sec:glwe-add-cipher}} GLWE Homomorphic Addition}}]
$\textsf{GLWE}_{S, \sigma}(\Delta M^{\langle 1 \rangle} + E^{\langle 1 \rangle} ) + \textsf{GLWE}_{S, \sigma}(\Delta M^{\langle 2 \rangle} + E^{\langle 2 \rangle}) $

$ = ( \{A_i^{\langle 1 \rangle}\}_{i=0}^{k-1}, \text{ } B^{\langle 1 \rangle}) + (\{A_i^{\langle 2 \rangle}\}_{i=0}^{k-1}, \text{ } B^{\langle 2 \rangle}) $

$ = ( \{A_i^{\langle 1 \rangle} + A_i^{\langle 2 \rangle}\}_{i=0}^{k-1}, \text{ } B^{\langle 1 \rangle} + B^{\langle 2 \rangle} ) $

$= \textsf{GLWE}_{S, \sigma}(\Delta(M^{\langle 1 \rangle} + M^{\langle 2 \rangle}) + E^{\langle 3 \rangle})$ \textcolor{red}{ $\rhd$ where $E^{\langle 3 \rangle} = E^{\langle 1  \rangle} + E^{\langle 2 \rangle}$}
\end{tcolorbox}


The phase identity is unconditional algebraically. Decoding returns $M^{\langle1\rangle}+M^{\langle2\rangle}\bmod t$ provided $\|E^{\langle1\rangle}+E^{\langle2\rangle}\|_\infty<\Delta/2$; for example, the sum of the two input noise bounds being less than $\Delta/2$ suffices.

$ $

%\noindent \textbf{\underline{Proof}}
\begin{myproof}
\begin{enumerate}
\item Define the following notations: \\
$A_0^{\langle 3 \rangle} = A_0^{\langle 1 \rangle} + A_0^{\langle 2 \rangle}$ \\
$A_1^{\langle 3 \rangle} = A_1^{\langle 1 \rangle} + A_1^{\langle 2 \rangle}$ \\
$\vdots$ \\
$A_{k-1}^{\langle 3 \rangle} = A_{k-1}^{\langle 1 \rangle} + A_{k-1}^{\langle 2 \rangle}$ \\
$E^{\langle 3 \rangle} = E^{\langle 1 \rangle} + E^{\langle 2 \rangle}$ \\
$B^{\langle 3 \rangle} = B^{\langle 1 \rangle} + B^{\langle 2 \rangle}$
\item Derive the following: \\
$B^{\langle 3 \rangle} = B^{\langle 1 \rangle} + B^{\langle 2 \rangle}$ \\
$ = \sum\limits_{i=0}^{k-1}{(A_i^{\langle 1 \rangle} \cdot S_i)} + \Delta \cdot M^{\langle 1 \rangle} + E^{\langle 1 \rangle} + \sum\limits_{i=0}^{k-1}{(A_i^{\langle 2 \rangle} \cdot S_i)} + \Delta \cdot M^{\langle 2 \rangle} + E^{\langle 2 \rangle}$ \\ 
$= \sum\limits_{i=0}^{k-1}{((A_i^{\langle 1 \rangle} + A_i^{\langle 2 \rangle}) \cdot S_i)} + \Delta \cdot (M^{\langle 1 \rangle} + M^{\langle 2 \rangle}) + (E^{\langle 1 \rangle} + E^{\langle 2 \rangle})$ \textcolor{red}{ $\rhd$ commutative and distributive rules} \\
$= \sum\limits_{i=0}^{k-1}{(A_i^{\langle 3 \rangle} \cdot S_i)} + \Delta \cdot (M^{\langle 1 \rangle} + M^{\langle 2 \rangle}) + E^{\langle 3 \rangle}$ \\
\item Since $B^{\langle 3 \rangle} = \sum\limits_{i=0}^{k-1}{(A_i^{\langle 3 \rangle} \cdot S_i)} + \Delta \cdot (M^{\langle 1 \rangle} + M^{\langle 2 \rangle}) + E^{\langle 3 \rangle}$, 

this means that $(A_0^{\langle 3 \rangle}, A_1^{\langle 3 \rangle}, A_2^{\langle 3 \rangle}, \ldots, A_{k-1}^{\langle 3 \rangle}, B^{\langle 3 \rangle})$ form the ciphertext: $\textsf{GLWE}_{S, \sigma}(\Delta \cdot (M^{\langle 1 \rangle} + M^{\langle 2 \rangle}) + E^{\langle 3 \rangle})$. \\
\item Thus, \\
$\textsf{GLWE}_{S, \sigma}(\Delta M^{\langle 1 \rangle} + E^{\langle 1 \rangle}) + \textsf{GLWE}_{S, \sigma}(\Delta M^{\langle 2 \rangle} + E^{\langle 2 \rangle})$ \\
$ = ( A_0^{\langle 1 \rangle} + A_0^{\langle 2 \rangle}, \text{ } A_1^{\langle 1 \rangle} + A_1^{\langle 2 \rangle},\ldots, A_{k-1}^{\langle 1 \rangle} + A_{k-1}^{\langle 2 \rangle}, \text{ } B^{\langle 1 \rangle} + B^{\langle 2 \rangle} )$ \\
$ = ( A_0^{\langle 3 \rangle}, A_1^{\langle 3 \rangle}, A_2^{\langle 3 \rangle}, \ldots,  A_{k-1}^{\langle 3 \rangle}, B^{\langle 3 \rangle})$ \\
$ = ( \{A_i^{\langle 3 \rangle}\}_{i=0}^{k-1}, B^{\langle 3 \rangle})$ \\
$= \textsf{GLWE}_{S, \sigma}(\Delta (M^{\langle 1 \rangle} + M^{\langle 2 \rangle}) + E^{\langle 3 \rangle})$

%\begin{flushright}
%\qedsymbol{} 
%\end{flushright}

\end{enumerate}
\end{myproof}

\subsection{Discussion}
\label{subsubsec:glwe-add-cipher-discuss} 

\para{Noise Elimination:} Secret-key subtraction removes the mask products exactly. Rounding removes the resulting error only while its coefficient magnitudes remain below $\Delta/2$. Valid inputs can add to an invalid output if the combined error exceeds that bound.

$ $

\para{Noise Growth:} The new error is $E^{\langle3\rangle}=E^{\langle1\rangle}+E^{\langle2\rangle}$. For independent zero-mean errors, variances add; with equal variance $\sigma^2$, summing $h$ independent samples gives variance $h\sigma^2$. Zero mean alone is insufficient. In general, $\operatorname{Var}(e_1+e_2)=\operatorname{Var}(e_1)+\operatorname{Var}(e_2)+2\operatorname{Cov}(e_1,e_2)$. Adding a ciphertext to itself produces $2e$ and variance $4\sigma^2$, not $2\sigma^2$. A deterministic coefficient bound uses the triangle inequality and requires no independence assumption.

$ $

\para{Plaintext Overflow when $t\nmid q$:} Mask products and ciphertext components are always interpreted modulo $q$. When $t\mid q$, plaintext wrap-around modulo $t$ is compatible with multiplication by $\Delta=q/t$. When $t\nmid q$ and $\Delta=\lfloor q/t\rfloor$, write each unreduced plaintext coefficient as $m=m'+vt$ with centered $m'$. The effective error is $e-(q\bmod t)v$. Under the assumptions of \autoref{subsubsec:lwe-noise-bound}, require $|e-(q\bmod t)v|<\Delta/2$. This condition depends on the overflow $v$; $q\gg t$ alone does not ensure it. The subsequent reduction modulo $t$ is part of message decryption.

 